% !TeX program = pdflatex
\documentclass[aspectratio=169,xcolor={dvipsnames,table}]{beamer}

\input{../../_en_preambulo_beamer}
% Comandos y macros estadísticas compartidas para las presentaciones Beamer en inglés (Metropolis)

% Conjuntos numéricos básicos
\newcommand{\R}{\mathbb{R}}
\newcommand{\N}{\mathbb{N}}
\newcommand{\Q}{\mathbb{Q}}
\newcommand{\Z}{\mathbb{Z}}
\newcommand{\set}[1]{\left\{ #1 \right\}}
\newcommand{\abs}[1]{\left|#1\right|}
\newcommand{\norm}[1]{\left\| #1 \right\|}

% Comandos estadísticos y probabilísticos del curso
\newcommand{\Var}{\operatorname{Var}}
\newcommand{\cov}{\operatorname{Cov}}
\newcommand{\corr}{\rho}
\newcommand{\comb}[2]{\begin{pmatrix} #1 \\ #2 \end{pmatrix}}
\newcommand{\s}{\sigma}
\newcommand{\particion}{\mathfrak{P}}
\newcommand{\card}[1]{\#\left( #1 \right)}
\newcommand{\seq}[3]{\set{#1}_{#2}^{#3}}
\newcommand{\E}{\mathbb{E}}
\newcommand{\Prob}{\mathbb{P}}

% Entornos pedagógicos y cajas en inglés académico natural (soportando tanto nombres en inglés como en español)
\newenvironment{discussion}{%
  \begin{alertblock}{Discussion Question}%
}{%
  \end{alertblock}%
}
\newenvironment{discusion}{%
  \begin{alertblock}{Discussion Question}%
}{%
  \end{alertblock}%
}

\newenvironment{reflection}{%
  \begin{alertblock}{Classroom Reflection}%
}{%
  \end{alertblock}%
}
\newenvironment{reflexion}{%
  \begin{alertblock}{Classroom Reflection}%
}{%
  \end{alertblock}%
}

\newenvironment{paradox}{%
  \begin{alertblock}{Exploratory Paradox}%
}{%
  \end{alertblock}%
}
\newenvironment{paradoja}{%
  \begin{alertblock}{Exploratory Paradox}%
}{%
  \end{alertblock}%
}

\newenvironment{motivation}{%
  \begin{exampleblock}{Motivation: Data Science in Practice}%
}{%
  \end{exampleblock}%
}
\newenvironment{motivacion}{%
  \begin{exampleblock}{Motivation: Data Science in Practice}%
}{%
  \end{exampleblock}%
}

\newenvironment{quickchallenge}{%
  \begin{block}{Quick Team Challenge}%
}{%
  \end{block}%
}
\newenvironment{retoclase}{%
  \begin{block}{Quick Team Challenge}%
}{%
  \end{block}%
}


\title{Bernoulli and Binomial Distributions}
\subtitle{Section 03.04 --- Dichotomous Processes and Combinatorial Inference}
\author[J. Castillo Colmenares]{Juliho Castillo Colmenares}
\institute[Tec de Monterrey]{Tecnológico de Monterrey}
\date{\vspace{-1.2cm}}

\begin{document}

% SLIDE 01: Title Page
\begin{frame}[plain]
	\titlepage
\end{frame}

% SLIDE 02: Roadmap
\begin{frame}{Roadmap --- Chapter 03: Discrete Random Variables}
	\begin{block}{Curricular Structure of Chapter 03}
		\small
		\begin{enumerate}
			\itemsep=0.04cm
			\item \textcolor{gray}{Section 03.01: Discrete Probability Mass Functions (PMF and Support)}
			\item \textcolor{gray}{Section 03.02: Cumulative Distribution Functions (Discrete CDF)}
			\item \textcolor{gray}{Section 03.03: Mathematical Expectation, Variance, and Moments}
			\item \textbf{\textcolor{TecRojo}{Section 03.04: Bernoulli and Binomial Distributions}} $\leftarrow$ \emph{Current focus}
			\item \textcolor{gray}{Section 03.05: Geometric and Negative Binomial Distributions}
			\item \textcolor{gray}{Section 03.06: Hypergeometric Distribution}
		\end{enumerate}
	\end{block}
	\vspace{-0.15cm}
	\begin{alertblock}{Curricular Objective of the Section}
		\footnotesize
		Master the algebraic foundations, moment properties, and combinatorial modeling of independent dichotomous trials, applying the Binomial distribution in quality control and systems engineering.
	\end{alertblock}
\end{frame}

% SLIDE 03: Motivation and Intuition
\begin{frame}{Motivation --- From Simple Dichotomy to Repeated Sampling}
	\begin{columns}[T]
		\begin{column}{0.48\textwidth}
			\begin{block}{The Fundamental Building Block}
				\small
				\begin{itemize}
					\itemsep=0.08cm
					\item In science and engineering, countless processes reduce to two incompatible states: **Success / Failure**, **Functional / Defective**, **1 / 0**.
					\item A single binary experiment is called a **Bernoulli Trial**.
					\item Modeling a single trial quantifies the baseline risk of an individual component.
				\end{itemize}
			\end{block}
		\end{column}
		\begin{column}{0.48\textwidth}
			\begin{alertblock}{The Need for Aggregation}
				\small
				\begin{itemize}
					\itemsep=0.08cm
					\item Engineers rarely assess isolated components: they inspect lots of size $n$, transmit $n$-bit packets, or conduct $n$ clinical trials.
					\item Repeating $n$ **independent** Bernoulli trials naturally gives rise to the **Binomial Distribution**, counting total accumulated successes.
				\end{itemize}
			\end{alertblock}
		\end{column}
	\end{columns}
\end{frame}

% SLIDE 04: The Bernoulli Trial
\begin{frame}{The Bernoulli Trial --- Formal Definition}
	\begin{block}{Bernoulli Random Variable ($X \sim \text{Bernoulli}(p)$)}
		\small
		A discrete random variable $X$ follows a Bernoulli distribution with success parameter $p \in (0, 1)$ if its support is $\mathcal{S}_X = \set{0, 1}$ and its probability mass function is given by:
		\begin{align*}
			f_X(x) = P(X = x) = p^x (1-p)^{1-x}, \quad \text{for } x \in \set{0, 1}.
		\end{align*}
		Where $q = 1-p$ denotes the complementary failure probability ($X=0$).
	\end{block}
	\vspace{-0.1cm}
	\begin{columns}[T]
		\begin{column}{0.48\textwidth}
			\begin{exampleblock}{Indicator Property}
				\footnotesize
				$X$ acts as the indicator function of an event $A \subset \Omega$: taking value $1$ if $\omega \in A$ and $0$ otherwise.
			\end{exampleblock}
		\end{column}
		\begin{column}{0.48\textwidth}
			\begin{alertblock}{Exact Normalization}
				\footnotesize
				Summing over the binary support verifies total probability: $f_X(0) + f_X(1) = (1-p) + p = 1$.
			\end{alertblock}
		\end{column}
	\end{columns}
\end{frame}

% SLIDE 05: Moments of the Bernoulli Trial
\begin{frame}{Moments of the Bernoulli Trial}
	\begin{block}{Derivation of First and Second Moments}
		\small
		Applying expectation and raw moment definitions over the dichotomous support $\set{0, 1}$:
		\begin{align*}
			\mathbb{E}[X] &= \sum_{x=0}^1 x \cdot P(X=x) = 0 \cdot (1-p) + 1 \cdot p = p, \\
			\mathbb{E}[X^2] &= \sum_{x=0}^1 x^2 \cdot P(X=x) = 0^2 \cdot (1-p) + 1^2 \cdot p = p.
		\end{align*}
	\end{block}
	\vspace{-0.15cm}
	\begin{columns}[T]
		\begin{column}{0.48\textwidth}
			\begin{exampleblock}{Bernoulli Variance}
				\footnotesize
				By the König-Huygens identity:
				\begin{align*}
					\text{Var}(X) = \mathbb{E}[X^2] - (\mathbb{E}[X])^2 = p - p^2 = pq.
				\end{align*}
			\end{exampleblock}
		\end{column}
		\begin{column}{0.48\textwidth}
			\begin{alertblock}{Maximum Uncertainty}
				\footnotesize
				The variance $pq = p(1-p)$ reaches its absolute maximum $\sigma^2 = 0.25$ precisely when $p = 0.50$ (maximum entropy).
			\end{alertblock}
		\end{column}
	\end{columns}
\end{frame}

% SLIDE 06: Genesis of the Binomial Distribution
\begin{frame}{Genesis of the Binomial Distribution --- Sum of i.i.d. Trials}
	\begin{block}{Additive Construction ($X \sim \text{Bin}(n, p)$)}
		\small
		Let $Y_1, Y_2, \dots, Y_n$ be independent and identically distributed (i.i.d.) random variables with $Y_i \sim \text{Bernoulli}(p)$. The variable counting total successes across all $n$ trials is defined as:
		\begin{align*}
			X = \sum_{i=1}^n Y_i.
		\end{align*}
	\end{block}
	\vspace{-0.15cm}
	\begin{columns}[T]
		\begin{column}{0.48\textwidth}
			\begin{exampleblock}{Sample Space and Support}
				\footnotesize
				\begin{itemize}
					\itemsep=0.04cm
					\item The sample space $\Omega$ contains $2^n$ ordered binary sequences.
					\item The support of $X$ is the finite discrete set $\mathcal{S}_X = \set{0, 1, 2, \dots, n}$.
				\end{itemize}
			\end{exampleblock}
		\end{column}
		\begin{column}{0.48\textwidth}
			\begin{alertblock}{Stochastic Independence}
				\footnotesize
				The joint probability of any specific sequence with $k$ successes and $n-k$ failures is exactly $p^k (1-p)^{n-k}$ due to mutual independence.
			\end{alertblock}
		\end{column}
	\end{columns}
\end{frame}

% SLIDE 07: Binomial Probability Mass Function
\begin{frame}{Binomial Probability Mass Function (PMF)}
	\begin{block}{Exact PMF Formula ($X \sim \text{Bin}(n, p)$)}
		\small
		For any integer $k \in \mathcal{S}_X = \set{0, 1, \dots, n}$, the number of distinct sequences containing exactly $k$ successes is the binomial coefficient $\comb{n}{k} = \dfrac{n!}{k!(n-k)!}$. Thus:
		\begin{align*}
			f_X(k) = P(X = k) = \comb{n}{k} p^k (1-p)^{n-k}.
		\end{align*}
	\end{block}
	\vspace{-0.15cm}
	\begin{columns}[T]
		\begin{column}{0.48\textwidth}
			\begin{exampleblock}{Binomial Theorem}
				\footnotesize
				Normalization is algebraically guaranteed by polynomial expansion:
				$\sum_{k=0}^n \comb{n}{k} p^k q^{n-k} = (p+q)^n = 1^n = 1$.
			\end{exampleblock}
		\end{column}
		\begin{column}{0.48\textwidth}
			\begin{alertblock}{Combinatorial Symmetry}
				\footnotesize
				If $X \sim \text{Bin}(n, p)$, the failure count $n-X$ follows an exact $\text{Bin}(n, 1-p)$ distribution.
			\end{alertblock}
		\end{column}
	\end{columns}
\end{frame}

% SLIDE 09: Moments of the Binomial Distribution
\begin{frame}{Moments of the Binomial Distribution}
	\begin{block}{Moment Theorem for $X \sim \text{Bin}(n, p)$}
		\small
		Since $X = \sum_{i=1}^n Y_i$ is the sum of $n$ independent Bernoulli variables with $\mathbb{E}[Y_i]=p$ and $\text{Var}(Y_i)=pq$, linearity of expectation and additivity of variance yield:
		\begin{align*}
			\mathbb{E}[X] &= \sum_{i=1}^n \mathbb{E}[Y_i] = np, \\
			\text{Var}(X) &= \sum_{i=1}^n \text{Var}(Y_i) = np(1-p) = npq.
		\end{align*}
	\end{block}
	\vspace{-0.15cm}
	\begin{columns}[T]
		\begin{column}{0.48\textwidth}
			\begin{exampleblock}{Standard Deviation ($\sigma_X$)}
				\footnotesize
				Absolute dispersion scales with the square root of sample size: $\sigma_X = \sqrt{npq}$.
			\end{exampleblock}
		\end{column}
		\begin{column}{0.48\textwidth}
			\begin{alertblock}{Relative Dispersion (CV)}
				\footnotesize
				The coefficient of variation shrinks as $n$ grows: $\text{CV} = \dfrac{\sigma_X}{\mu_X} = \sqrt{\dfrac{q}{np}} \propto \dfrac{1}{\sqrt{n}}$.
			\end{alertblock}
		\end{column}
	\end{columns}
\end{frame}

% SLIDE 10: Shape of the Distribution, Skewness, and Mode
\begin{frame}{Distribution Shape, Skewness, and Binomial Mode}
	\begin{columns}[T]
		\begin{column}{0.48\textwidth}
			\begin{block}{Binomial Mode Theorem}
				\small
				The integer $k^* \in \mathcal{S}_X$ that maximizes the probability $f_X(k)$ satisfies:
				\begin{align*}
					k^* = \lfloor (n+1)p \rfloor.
				\end{align*}
				If $(n+1)p \in \mathbb{Z}^+$, the distribution is **bimodal** across $k^* - 1$ and $k^*$.
			\end{block}
		\end{column}
		\begin{column}{0.48\textwidth}
			\begin{alertblock}{Skewness ($\gamma_1$)}
				\small
				The standardized third central moment measures asymmetry around the mean:
				\begin{align*}
					\gamma_1 = \dfrac{q - p}{\sqrt{npq}} = \dfrac{1 - 2p}{\sqrt{np(1-p)}}.
				\end{align*}
				It is perfectly symmetric ($\gamma_1=0$) iff $p=0.5$.
			\end{alertblock}
		\end{column}
	\end{columns}
\end{frame}

% SLIDE 17: Python Lab --- Vectorized PMF and CDF Validation
\begin{frame}[fragile]{Python Lab --- Vectorized PMF and CDF Validation (`numpy`)}
	\vspace{-0.15cm}
	\begin{block}{Single Source of Truth Implementation (\texttt{03.04\_bernoulli\_binomial.py})}
		\lstinputlisting[language=Python, firstline=13, lastline=34, basicstyle=\fontsize{5.8pt}{6.8pt}\ttfamily]{../../code/03_variables_aleatorias_discretas/03.04_bernoulli_binomial.py}
	\end{block}
\end{frame}

% SLIDE 18: Python Lab --- Monte Carlo Simulation and Moments
\begin{frame}[fragile]{Python Lab --- Monte Carlo Simulation and Moments (`numpy`)}
	\begin{block}{Large-Scale Simulation ($N=100,000$ Bernoulli Trials)}
		\lstinputlisting[language=Python, firstline=39, lastline=56, basicstyle=\fontsize{6.3pt}{7.5pt}\ttfamily]{../../code/03_variables_aleatorias_discretas/03.04_bernoulli_binomial.py}
	\end{block}
\end{frame}

% SLIDE 19: Python Lab --- Mode and Skewness Analysis
\begin{frame}[fragile]{Python Lab --- Mode and Skewness Analysis (`scipy.stats`)}
	\begin{block}{Computational Verification of Peak and Asymmetry}
		\lstinputlisting[language=Python, firstline=61, lastline=75, basicstyle=\fontsize{6.3pt}{7.5pt}\ttfamily]{../../code/03_variables_aleatorias_discretas/03.04_bernoulli_binomial.py}
	\end{block}
\end{frame}


% SLIDE 20: Python Lab --- Verified Console Output
\begin{frame}[fragile]{Python Lab --- Verified Console Output}
	\begin{block}{Execution console (\texttt{python 03.04\_bernoulli\_binomial.py})}
		\fontsize{6pt}{7pt}\selectfont\ttfamily
=== Section 03.04: Bernoulli and Binomial Lab ===\\[0.08cm]
{[1]} Vectorized PMF and CDF validation (n=15, p=0.2):\\
\ \ \ Total support mass P(X<=n): 1.000000 (target: 1.0)\\
\ \ \ Maximum PMF error versus SciPy: 6.66e-16\\
\ \ \ Maximum CDF error versus SciPy: 2.78e-15\\[0.08cm]
{[2]} Monte Carlo LLN simulation (N=100,000):\\
\ \ \ Empirical mean: 2.9970 (target: 3.0000, error: 0.003020)\\
\ \ \ Empirical variance: 2.3827 (target: 2.4000, error: 0.017289)\\[0.08cm]
{[3]} Mode and skewness analysis:\\
\ \ \ Computational mode: 3 (theoretical floor((n+1)p): 3)\\
\ \ \ Exact skewness: 0.387298 (SciPy target: 0.387298)\\[0.08cm]
=== All numerical checks passed exact tolerances. ===
	\end{block}
\end{frame}

% SLIDE 20: Conclusion & Stochastic Synthesis
\begin{frame}{Conclusion --- Stochastic Synthesis and Didactic Bridge}
	\begin{columns}[T]
		\begin{column}{0.48\textwidth}
			\begin{block}{Conceptual Summary}
				\small
				\begin{itemize}
					\itemsep=0.08cm
					\item **Bernoulli ($p$):** Basic binary building block, $\mu = p$, $\sigma^2 = pq$.
					\item **Binomial ($n, p$):** Sum of $n$ independent trials, $\mu = np$, $\sigma^2 = npq$.
					\item **Invariance:** Python Monte Carlo simulations verified theoretical moments with $< 10^{-4}$ error.
				\end{itemize}
			\end{block}
		\end{column}
		\begin{column}{0.48\textwidth}
			\begin{alertblock}{Bridge to Section 03.05}
				\small
				What happens if instead of fixing total trials $n$ and counting successes, **we fix the desired number of successes $r$ and count the required trials**?
				
				$\to$ This leads directly to the **Geometric Distribution** ($r=1$) and the **Negative Binomial Distribution** ($r \ge 1$).
			\end{alertblock}
		\end{column}
	\end{columns}
\end{frame}

\end{document}
